
% peaks: m/z 91,105,122,150,178  heights 35,100,62,18,48
\newcommand{\stemUP}[3]{% x, height(0..100), color
  \draw[#3,fig xbold] (#1,0) -- (#1,{#2/100*2.6});}
\newcommand{\stemDN}[3]{% x, height(0..100), color
  \draw[#3,fig xbold] (#1,0) -- (#1,{-#2/100*2.6});}

\begin{tikzpicture}[fig]

% =====================================================================
% PANEL A — aligned (shift 0): perfect match
% =====================================================================
\begin{scope}[shift={(0,0)}]
  \node[anchor=west,panel title] at (0,3.5)
    {(a) Aligned spectra \;—\; identical, $\theta=0$};
  % axis: x from 80..190 mapped to 0..9 ; use raw m/z with xscale
  \begin{scope}[xscale=9/110, xshift=-80*9/110 cm]
    % zero baseline
    \draw[black!41,fig line] (80,0) -- (192,0);
    \draw[-{Latex[length=2mm]},black!41] (80,0) -- (194,0);
    % A up (red), B down (blue) at same positions
    \stemUP{91}{35}{BrickRed}\stemDN{91}{35}{RoyalBlue}
    \stemUP{105}{100}{BrickRed}\stemDN{105}{100}{RoyalBlue}
    \stemUP{122}{62}{BrickRed}\stemDN{122}{62}{RoyalBlue}
    \stemUP{150}{18}{BrickRed}\stemDN{150}{18}{RoyalBlue}
    \stemUP{178}{48}{BrickRed}\stemDN{178}{48}{RoyalBlue}
    % m/z ticks
    \foreach \m in {91,105,122,150,178}{
      \node[black!41,font=\scriptsize] at (\m,-3.05) {\m};}
    \node[black!41,font=\scriptsize] at (193,-0.35) {$m/z$};
  \end{scope}
  \node[BrickRed,font=\small] at (8.4,2.7) {A};
  \node[RoyalBlue,font=\small] at (8.4,-2.7) {B};
  \node[anchor=west,font=\small] at (0.1,-3.55)
    {cosine $=\mathbf{1.00}$ \qquad Wasserstein $=\mathbf{0.0}$};
\end{scope}

% =====================================================================
% PANEL B — shifted (B moved +8 Da): peaks no longer overlap
% =====================================================================
\begin{scope}[shift={(0,-8.4)}]
  \node[anchor=west,panel title] at (0,3.5)
    {(b) B shifted $+8$\,Da \;—\; peaks miss, $\theta\!\to\!90^\circ$};
  \begin{scope}[xscale=9/110, xshift=-80*9/110 cm]
    \draw[black!41,fig line] (80,0) -- (192,0);
    \draw[-{Latex[length=2mm]},black!41] (80,0) -- (194,0);
    % A up at original positions
    \stemUP{91}{35}{BrickRed}\stemUP{105}{100}{BrickRed}\stemUP{122}{62}{BrickRed}
    \stemUP{150}{18}{BrickRed}\stemUP{178}{48}{BrickRed}
    % B down at +8
    \stemDN{99}{35}{RoyalBlue}\stemDN{113}{100}{RoyalBlue}\stemDN{130}{62}{RoyalBlue}
    \stemDN{158}{18}{RoyalBlue}\stemDN{186}{48}{RoyalBlue}
    % transport arrows (Wasserstein: move each blue back onto red), amber
    \foreach \bx/\ax in {99/91,113/105,130/122,158/150,186/178}{
      \draw[yellow,-{Latex[length=1.6mm]},fig line,opacity=0.8]
        (\bx,-0.5) -- (\ax,-0.5);}
    \node[yellow,font=\scriptsize] at (135,-1.15) {transport distance $=8$ each};
    \foreach \m in {91,105,122,150,178}{
      \node[BrickRed,font=\scriptsize] at (\m,2.95) {\m};}
    \node[black!41,font=\scriptsize] at (193,-0.35) {$m/z$};
  \end{scope}
  \node[BrickRed,font=\small] at (8.6,2.4) {A};
  \node[RoyalBlue,font=\small] at (8.6,-2.4) {B};
  \node[anchor=west,font=\small] at (0.1,-3.55)
    {cosine $=\mathbf{0.00}$ (no shared peaks) \qquad Wasserstein $=\mathbf{21.1}$};
\end{scope}

% =====================================================================
% PANEL C — response curves vs shift
% =====================================================================
\begin{scope}[shift={(0,-16.0)}]
  \node[anchor=west,panel title] at (0,0.55)
    {(c) Response vs.\ how far B is shifted};
  \begin{axis}[
      at={(0,0cm)}, anchor=north west,
      width=10.6cm, height=6.2cm,
      xlabel={$|\text{shift of }B|$ (Da)},
      xlabel style={font=\small,color=black!41},
      axis y line*=left, axis x line*=bottom,
      xmin=0, xmax=20, ymin=0, ymax=1.05,
      ylabel={cosine similarity},
      ylabel style={font=\small,color=RoyalBlue},
      ytick={0,0.25,0.5,0.75,1.0},
      yticklabel style={color=RoyalBlue,font=\scriptsize},
      xtick={0,4,8,12,16,20},
      xticklabel style={color=black!41,font=\scriptsize},
      tick style={color=black!41},
      every axis plot/.append style={fig heavy},
    ]
    % cosine: 1 at 0, then 0 (with spurious spikes at 14,17)
    \addplot[RoyalBlue,mark=*,mark size=1.2pt] coordinates {
      (0,1.0)(1,0)(2,0)(3,0)(4,0)(5,0)(6,0)(7,0)(8,0)(9,0)(10,0)
      (11,0)(12,0)(13,0)(14,0.198)(15,0)(16,0)(17,0.350)(18,0)(19,0)(20,0)};
    \node[RoyalBlue,font=\scriptsize,anchor=south west] at (axis cs:14,0.20) {spurious};
    \node[RoyalBlue,font=\scriptsize,anchor=south west] at (axis cs:14,0.05) {re-alignment};
  \end{axis}
  \begin{axis}[
      at={(0,0cm)}, anchor=north west,
      width=10.6cm, height=6.2cm,
      axis y line*=right, axis x line=none,
      xmin=0, xmax=20, ymin=0, ymax=56,
      ylabel={Wasserstein distance},
      ylabel style={font=\small,color=yellow},
      ytick={0,10,20,30,40,50},
      yticklabel style={color=yellow,font=\scriptsize},
      tick style={color=yellow},
    ]
    % Wasserstein: linear 2.63 * shift
    \addplot[yellow,fig heavy,dashed] coordinates {
      (0,0)(20,52.6)};
    \addplot[yellow,only marks,mark=*,mark size=1.1pt] coordinates {
      (0,0)(2,5.25)(4,10.5)(6,15.8)(8,21.1)(10,26.3)(12,31.5)
      (14,36.8)(16,42.1)(18,47.3)(20,52.6)};
  \end{axis}
  % legend
  \node[anchor=north west,font=\scriptsize] at (7.0,0.5)
    {\textcolor{RoyalBlue}{\rule[0.4ex]{10pt}{1.3pt}}~cosine};
  \node[anchor=north west,font=\scriptsize] at (7.0,0.1)
    {\textcolor{yellow}{\rule[0.4ex]{10pt}{1.3pt}}~Wasserstein};
\end{scope}

% =====================================================================
% PANEL D — response to scaling B's intensity (peaks aligned)
% =====================================================================
\begin{scope}[shift={(0,-24.4)}]
  \node[anchor=west,panel title] at (0,0.55)
    {(d) Response vs.\ B's peak intensity \;(positions aligned)};
  \begin{axis}[
      at={(0,0cm)}, anchor=north west,
      width=10.6cm, height=6.2cm,
      xlabel={intensity scale of B ($\times$)},
      xlabel style={font=\small,color=black!41},
      axis y line*=left, axis x line*=bottom,
      xmin=1, xmax=3, ymin=0, ymax=1.18,
      ylabel={cosine similarity},
      ylabel style={font=\small,color=RoyalBlue},
      ytick={0,0.25,0.5,0.75,1.0},
      yticklabel style={color=RoyalBlue,font=\scriptsize},
      xtick={1,1.5,2,2.5,3},
      xticklabel style={color=black!41,font=\scriptsize},
      tick style={color=black!41},
    ]
    % cosine: flat at 1.0 regardless of intensity scale
    \addplot[RoyalBlue,fig heavy,mark=*,mark size=1.2pt] coordinates {
      (1,1)(1.25,1)(1.5,1)(1.75,1)(2,1)(2.5,1)(3,1)};
    \node[RoyalBlue,font=\scriptsize,anchor=south] at (axis cs:2,1.02)
      {cosine unchanged $=1.00$ (scale-invariant)};
  \end{axis}
  \begin{axis}[
      at={(0,0cm)}, anchor=north west,
      width=10.6cm, height=6.2cm,
      axis y line*=right, axis x line=none,
      xmin=1, xmax=3, ymin=0, ymax=340,
      ylabel={Wasserstein distance},
      ylabel style={font=\small,color=yellow},
      ytick={0,100,200,300},
      yticklabel style={color=yellow,font=\scriptsize},
      tick style={color=yellow},
    ]
    % Wasserstein grows linearly with added mass
    \addplot[yellow,fig heavy,dashed] coordinates {(1,0)(3,312.8)};
    \addplot[yellow,only marks,mark=*,mark size=1.1pt] coordinates {
      (1,0)(1.25,39.1)(1.5,78.2)(1.75,117.3)(2,156.4)(2.5,234.6)(3,312.8)};
    \node[yellow,font=\scriptsize,anchor=west] at (axis cs:1.05,250)
      {Wasserstein rises};
    \node[yellow,font=\scriptsize,anchor=west] at (axis cs:1.05,222)
      {with added mass};
  \end{axis}
\end{scope}

\end{tikzpicture}